\section{Introduction}
\label{sec:intro}

A growing share of agent work is optimization: a number that should go down (runtime, latency, loss) or up
(accuracy, a benchmark score), with tests that must keep passing. Karpathy's \emph{autoresearch}~\cite{autoresearch}
made a simple loop popular for this: change one thing, run the measurement, keep the change if it wins, discard
it if not, and log every round. Many variants now ship as prompts or agent skills. They are written as if the
loop obviously helps. We did not find one that tests that.

The question has a practical edge. An agent told to ``make the suite faster'' can do it by testing less. An agent
told to ``raise accuracy'' can tune on its validation split until the gain does not survive a held-out set. An
agent reading a noisy benchmark can mistake noise for progress. A loop recipe claims to guard against all three.
Whether it does is an empirical question about agents, not about the recipe's wording.

We packaged the loop as an agent skill, \emph{hill-climb}, with a small helper script that freezes the setup,
runs gates, measures repeatedly, decides keep or revert against a noise band, and logs each round. We then built a
benchmark to test it (Section~\ref{sec:bench}): four optimization tasks, each with an objective score and a
hidden gate the agent never sees, run in isolated \texttt{claude -p} sessions and scored from outside. We compare
the skill against two controls on the same model: the plain task, and the task plus an instruction to keep
iterating.

Our findings, in order of how much we trust them:
\begin{enumerate}
\item \textbf{The skill prevented a weakened test suite on Opus.} Asked to speed up a test suite, 5 of 8 Opus runs
  without the skill shrank a slow fixture until the suite stopped killing a hidden mutant, and said the suite was no
  weaker. With the skill, 0 of 4 did (Section~\ref{sec:results}). The effect rests on one mutant and small samples.
\item \textbf{It changed how agents work, not how far they got on easy tasks.} With the revised skill, Sonnet agents
  logged a median of 7 measured rounds per run (13.5 on Opus); with the first version, 3, and 7 of 24 runs logged
  one round or none. On the speed task the first version mostly stopped after one large rewrite. On three
  of four tasks the final scores tie a plain loop, at roughly twice the cost.
\item \textbf{On a noisy benchmark it still trails.} On a tuning task with 8\% run-to-run noise, a free-form loop
  reached a 1.65$\times$ improvement and the best skill version 1.43$\times$.
\item \textbf{The benchmark debugged the skill.} The noisy task exposed two bugs in the helper, both of which made
  correct changes look like noise. Fixing them lifted the skill from 1.31$\times$ to 1.43$\times$ there
  (Section~\ref{sec:fix}).
\item \textbf{Blind judge panels are consistent, not correct.} On 100 real headline A/B tests, three-judge panels
  pick the winner 65 to 69\% of the time however much they agree (Section~\ref{sec:taste}).
\end{enumerate}

No Sonnet run shipped a broken gate, and only one Sonnet run claimed a gain the outside measurement did not
support. The failure that mattered was quieter: an honest-sounding claim, on Opus, that a cheaper test suite was
just as good.
